\scanpage{18}
% lines 1–4
exprimant que un certain \uline{autom} de $E_{\widetilde S}$ ($\widetilde S={}$ schém. des nbs duaux) défini par $\rho(x)$ invarie $\varphi_{\widetilde S}$ -- donc invarie $[\ ,\ ]_{\varphi_{\widetilde S}}$ -- et en interprétant (53) de façon analogue.

% lines 5–7
\uline{Construction.} L'algèbre de Lie \deleted{$\unsure{\mathfrak g_E}$} $\unsure{\mathfrak g|g}=E'$ opérant sur le module $E'$ (par représentation adjointe

% lines 8
\[
\rho_{E'}(\deleted{\unsure{\eta'|x}}\xi')x'\overset{\mathrm{dfn}}{=}[\eta',x']_{\varphi\unsure{*|'}}\,, \tag*{(57)}
\]

% lines 9–10
on la fait opérer sur $E=\check E'$ par contragrédiente,

% lines 11
\[
\rho_E(\xi')=-{}^t\rho_{E'}\deleted{\unsure{E}}(\eta') \tag*{(58)}
\]

% lines 12
i.e. par exigeant

% lines 13
\[
\langle\rho_E(\eta')x,x'\rangle+\langle x,\rho_{E'}(\eta')x'\rangle=0 \tag*{(59)}
\]

% lines 14
i.e.

% lines 15–17
\[
\begin{aligned}
\langle\rho_E(\eta')x,x'\rangle
&=-\langle x,\underbrace{[\eta',x']_{\varphi *}}_{f\alpha'(\eta'\wedge x')}\rangle\\
&\unsure{-|=}\\[-.5ex]
&-\varphi(\deleted{\unsure{x,}}\alpha'(\overset{\Vert}{\eta'}\wedge x'),x)
\end{aligned}
\tag*{(59 bis)}
\]

% lines 18–20
Le fait que $\mathfrak g$ opère sur $E'$ (de façon compatible avec les crochets) implique qu'il opère sur $E=\check E'$, i.e. qu'on a (55).

% lines 21
La formule (56) s'écrit \deleted{en prenant le}

% lines foot
% archive sheet number: 17
